\begin{theorem}[Sector scalars from a product tensor identity]
    \label{thm:periodic_product_tensor_phase_extraction}
    \lean{TNLean.Algebra.PiTensorProductPhase.exists_kappa_of_piTensorProduct_eq_smul}
    \leanok
    Let $z\in\C$ be nonzero. For each sector $v=0,\ldots,m-1$, let
    $\{A_v^i\}_i$ and $\{B_v^i\}_i$ be families of complex matrices with the
    same rectangular size. Suppose that for every choice of letters
    $\sigma=(\sigma_v)_v$,
    \begin{align}
        \bigotimes_{v=0}^{m-1} A_v^{\sigma_v}
        &=
        z\,\bigotimes_{v=0}^{m-1} B_v^{\sigma_v}.
        \notag
    \end{align}
    Suppose also that for each $v$ there is a reference letter $i_v$ and a
    matrix entry $(r_v,c_v)$ with
    $(B_v^{i_v})_{r_v,c_v}\ne 0$. Then there are scalars
    $\kappa_v\in\C$ such that
    \begin{align}
        \prod_{v=0}^{m-1}\kappa_v&=z,
        \notag\\
        A_v^i&=\kappa_v B_v^i
        \notag
    \end{align}
    for every sector $v$ and every letter $i$.

    This is the scalar extraction used in
    \cite[Appendix~A, lines~1072--1080]{DeLasCuevas2017Irreducible}, after the
    contraction with the maps $F_u$ and $\Omega_u$ has produced the displayed
    product-tensor identity.
\end{theorem}

\begin{proof}\leanok
    For each sector choose the coordinate functional at the nonzero reference
    entry and normalize it so that it takes the value $1$ on
    $B_v^{i_v}$. Applying the tensor product of these functionals to the
    reference tensor identity gives $\prod_v\kappa_v=z$. To prove the
    proportionality in sector $v$, replace only the $v$-th functional by the
    coordinate functional at an arbitrary entry $(r,c)$ and apply the tensor
    identity with $\sigma_v=i$ and all other letters fixed at their references.
    The product over the other sectors is nonzero because its product with
    $\kappa_v$ is $z$, so cancellation gives
    $(A_v^i)_{r,c}=(\kappa_v B_v^i)_{r,c}$ for every entry.
\end{proof}

\begin{theorem}[Product-one sector scalars after choosing a root]
    \label{thm:periodic_product_one_phase_extraction}
    \lean{TNLean.Algebra.PiTensorProductPhase.exists_kappa_product_one_of_piTensorProduct_eq_root_smul}
    \leanok
    \uses{thm:periodic_product_tensor_phase_extraction}
    Let $m>0$, let $z\in\C$ be nonzero, and choose $\xi\in\C$ with
    $\xi^m=z$. For each sector $v=0,\ldots,m-1$, let $\{A_v^i\}_i$ and
    $\{B_v^i\}_i$ be families of complex matrices with the same rectangular
    size. Suppose that for every choice of letters $\sigma=(\sigma_v)_v$,
    \begin{align}
        \bigotimes_{v=0}^{m-1} A_v^{\sigma_v}
        &=
        z\,\bigotimes_{v=0}^{m-1} B_v^{\sigma_v}.
        \notag
    \end{align}
    Suppose also that each sector has a reference letter $i_v$ and a matrix
    entry $(r_v,c_v)$ with $(B_v^{i_v})_{r_v,c_v}\ne0$. Then there are scalars
    $\kappa_v\in\C$ such that
    \begin{align}
        \prod_{v=0}^{m-1}\kappa_v&=1,
        \notag\\
        A_v^i&=\kappa_v\xi\,B_v^i
        \notag
    \end{align}
    for every sector $v$ and every letter $i$.

    This is the normalization in
    \cite[Appendix~A, lines~1072--1080]{DeLasCuevas2017Irreducible}: after the
    preceding translation step has made the scalar independent of the sector
    label, choosing a common $m$-th root of that scalar leaves a product-one
    family of sector scalars.
\end{theorem}

\begin{proof}\leanok
    Apply Theorem~\ref{thm:periodic_product_tensor_phase_extraction} to obtain
    scalars $\mu_v$ with $\prod_v\mu_v=z$ and $A_v^i=\mu_vB_v^i$. Since
    $m>0$, the identities $\xi^m=z$ and $z\ne0$ imply $\xi\ne0$. Put
    $\kappa_v=\mu_v/\xi$. Then
    $\prod_v\kappa_v=(\prod_v\mu_v)/\xi^m=z/z=1$, and
    $A_v^i=(\kappa_v\xi)B_v^i$.
\end{proof}

\begin{theorem}[Left-canonical scalar normalization]
    \label{thm:left_canonical_scalar_normalization}
    \lean{MPSTensor.norm_eq_one_of_leftCanonical_smul}
    \leanok
    Let $A$ have positive bond dimension, and let $c\in\C$. If both $A$ and
    the scalar multiple $cA$ are left-canonical, then $|c|=1$.
\end{theorem}

\begin{proof}\leanok
    The two left-canonical equations give
    \begin{align}
        \sum_i (A^i)^\dagger A^i&=\Id,
        \notag\\
        \sum_i (cA^i)^\dagger(cA^i)&=\Id.
        \notag
    \end{align}
    Substituting the first equation into the second gives
    $(\bar c c)\Id=\Id$. Since the bond dimension is positive, comparison of
    any diagonal entry gives $\bar c c=1$, hence $|c|=1$.
\end{proof}

\begin{theorem}[Left-canonical sector scalar normalization]
    \label{thm:left_canonical_sector_scalar_normalization}
    \lean{MPSTensor.kappa_norm_eq_one_of_leftCanonical_smul}
    \leanok
    \uses{thm:left_canonical_scalar_normalization}
    Let $A$ and $B$ be left-canonical tensors with the same positive bond
    dimension. Suppose that $|\xi|=1$ and that
    $A^i=(\kappa\xi)B^i$ for every physical index~$i$. Then $|\kappa|=1$.

    This is the scalar normalization in
    \cite[Appendix~A, lines~1082--1084]{DeLasCuevas2017Irreducible}, after the
    sectorwise proportionality relation has been obtained.
\end{theorem}

\begin{proof}\leanok
    By Theorem~\ref{thm:left_canonical_scalar_normalization}, applied to $B$
    and the scalar $\kappa\xi$, one has $|\kappa\xi|=1$. Since $|\xi|=1$, the
    multiplicativity of the complex norm gives $|\kappa|=1$.
\end{proof}

\begin{definition}[Off-diagonal sector decomposition]%
    \label{def:cyclic_sector_decomp}
    \lean{MPSTensor.IsCyclicSectorDecomp, MPSTensor.IsCyclicSectorDecompWith}
    \leanok
    \uses{def:periodic_tensor}
    A family of compressed sector tensors
    $(C_k)_{k=0}^{m-1}$ on corner bond spaces forms an \emph{off-diagonal
        decomposition} of the blocked tensor $A^{(m)}$ if there exist orthogonal
    projections $P_0,\ldots,P_{m-1}$ on the ambient bond space such that
    $\sum_k P_k = \Id$, the projections satisfy the cyclic-shift relation
    $\E_A^\dagger(P_{k+1}) = P_k$ with the index $k+1$ taken cyclically
    modulo~$m$, each $P_k$ commutes with every blocked letter
    $P_k A^{(m)}_i = A^{(m)}_i P_k$, and
    $V^{(N)}(C_k)_\sigma = \tr(P_k (A^{(m)})^\sigma)$ for every word~$\sigma$,
    and each block carries a $\ast$-algebra isomorphism
    $\varphi_k : M_{\dim C_k}(\mathbb{C}) \xrightarrow{\sim}
        P_k \cdot M_D(\mathbb{C}) \cdot P_k$
    (multiplicative and adjoint-preserving) intertwining the compressed adjoint
    transfer map with the sector adjoint transfer map on the corner of~$P_k$.
    We shall sometimes keep these projections and corner isomorphisms as
    specified witnesses; forgetting the witnesses gives the displayed
    decomposition.
    This is the inverse cyclic indexing of the off-diagonal convention in
    \cite[Appendix~A]{DeLasCuevas2017Irreducible}. There the same adjoint
    transfer map is written $E^*$, and the source projectors satisfy
    $E^*(P_u)=P_{u+1}$ together with
    $A^i=\sum_u P_u A^i P_{u+1}$. Here $P_k$ corresponds to the source
    sector $u=-k \pmod m$, which turns the source shift into the displayed
    adjoint-transfer relation.
\end{definition}

\begin{theorem}[Single-site cyclic grading]
    \label{thm:offdiag_shift_adjoint_cyclic_shift}
    \lean{MPSTensor.offDiag_shift_of_adjoint_cyclic_shift}
    \leanok
    \uses{def:transfer_map}
    Let $A$ be left-canonical, and let $P_0,\ldots,P_{m-1}$ be orthogonal
    projections on the bond space. If the adjoint transfer map satisfies
    $\E_A^\dagger(P_{k+1})=P_k$ for all $k$, with indices taken modulo $m$,
    then $P_{k+1}A^i=A^iP_k$ for every physical letter $i$.
    This is the single-site grading behind
    \cite[eq.~Aoffdiag]{DeLasCuevas2017Irreducible}; the compressed blocks of
    \cite[eq.~Auprop]{DeLasCuevas2017Irreducible} use the inverse cyclic
    indexing convention of Definition~\ref{def:cyclic_sector_decomp}.
\end{theorem}

\begin{proof}\leanok
    Put $K_i=(A^i)^\dagger$ and let $\Phi_K$ be the Kraus map of the family
    $K$. Since $A$ is left-canonical, $\Phi_K$ is unital. The projection
    identities and the shift relation give
    \begin{align}
        \Phi_K(P_{k+1}^\dagger P_{k+1})
        &=
        \Phi_K(P_{k+1})^\dagger\Phi_K(P_{k+1}),
        \notag
    \end{align}
    because both sides are equal to $P_k$. The equality case in the
    Kadison--Schwarz inequality therefore gives
    $P_{k+1}K_i^\dagger=K_i^\dagger\Phi_K(P_{k+1})$. Substituting
    $K_i^\dagger=A^i$ and $\Phi_K(P_{k+1})=P_k$ gives
    $P_{k+1}A^i=A^iP_k$.
\end{proof}

\begin{theorem}[Off-diagonal reconstruction]
    \label{thm:eq_sum_offdiag_adjoint_cyclic_shift}
    \lean{MPSTensor.eq_sum_offDiag_of_adjoint_cyclic_shift}
    \leanok
    \uses{thm:offdiag_shift_adjoint_cyclic_shift}
    Under the hypotheses of Theorem~\ref{thm:offdiag_shift_adjoint_cyclic_shift},
    assume in addition that $\sum_{u=0}^{m-1} P_u=\Id$. Then every physical
    letter decomposes as $A^i=\sum_{u=0}^{m-1}P_{u+1}A^iP_u$.
    This is \cite[eq.~Aoffdiag]{DeLasCuevas2017Irreducible}, written in the
    same inverse cyclic indexing convention.
\end{theorem}

\begin{proof}\leanok
    Insert $\sum_u P_u=\Id$ on the right of $A^i$, then use
    Theorem~\ref{thm:offdiag_shift_adjoint_cyclic_shift} and
    $P_u^2=P_u$ term by term.
\end{proof}

\begin{theorem}[Word cyclic grading]
    \label{thm:offdiag_shift_eval_word_adjoint_cyclic_shift}
    \lean{MPSTensor.offDiag_shift_evalWord_of_adjoint_cyclic_shift}
    \leanok
    \uses{thm:offdiag_shift_adjoint_cyclic_shift}
    Under the hypotheses of Theorem~\ref{thm:offdiag_shift_adjoint_cyclic_shift},
    a word $w=i_1\cdots i_\ell$ satisfies $P_{k+\ell} A^w=A^w P_k$.
    Thus the single-site grading accumulates one cyclic step for each letter of
    the word.
\end{theorem}

\begin{proof}\leanok
    Induct on the length of $w$, using
    Theorem~\ref{thm:offdiag_shift_adjoint_cyclic_shift} for the first letter
    and the induction hypothesis for the remaining word.
\end{proof}

\begin{definition}[Corner-transition letters]
    \label{def:periodic_corner_letter}
    \lean{MPSTensor.cornerLetter}
    \leanok
    \uses{def:periodic_tensor}
    Let $P_0,\ldots,P_{m-1}$ be cyclic sector projections for a tensor $A$.
    The one-site corner-transition letters of
    \cite[Appendix~A]{DeLasCuevas2017Irreducible} are
    $A_u^i=P_u A^i P_{u+1}$, with the index $u+1$ taken cyclically
    modulo~$m$.
\end{definition}

\begin{definition}[Repeated corner products]
    \label{def:periodic_corner_product}
    \lean{MPSTensor.cornerProd}
    \leanok
    \uses{def:periodic_corner_letter}
    For a word $\mathbf i=(i_1,\ldots,i_\ell)$, define
    \begin{align}
        F_u^{\mathbf i}
        &=
        A_u^{i_1}A_{u+1}^{i_2}\cdots A_{u+\ell-1}^{i_\ell}.
        \notag
    \end{align}
    The empty word gives $F_u^\varnothing=P_u$, and a one-letter word gives
    $F_u^{(i)}=A_u^i$. This is the notation of
    \cite[Appendix~A]{DeLasCuevas2017Irreducible}; in the Case~3 contraction
    the length is $\ell=mN_0$.
\end{definition}

\begin{lemma}[Concatenation of corner products]
    \label{lem:periodic_corner_product_append}
    \lean{MPSTensor.cornerProd_append}
    \leanok
    \uses{def:periodic_corner_product}
    If the sector maps $P_u$ are orthogonal projections, then for any two
    words $\mathbf i$ and $\mathbf j$,
    $F_u^{\mathbf i\mathbf j}
    =F_u^{\mathbf i}F_{u+|\mathbf i|}^{\mathbf j}$, where indices are taken
    cyclically modulo~$m$.
\end{lemma}

\begin{proof}\leanok
    \uses{def:periodic_corner_product}
    Induct on the word $\mathbf i$. For the empty word the identity is
    $P_uF_u^{\mathbf j}=F_u^{\mathbf j}$, which follows from the idempotent
    part of the orthogonal-projection hypothesis. For a nonempty word, remove
    the first letter and apply the induction hypothesis at the next cyclic
    sector.
\end{proof}

\begin{lemma}[Repeated-block matrix identity]
    \label{lem:periodic_corner_product_blockmatch_pow}
    \lean{MPSTensor.cornerProd_blockMatch_pow}
    \leanok
    \uses{def:periodic_corner_product, lem:periodic_corner_product_append}
    Let $P_0,\ldots,P_{m-1}$ and $Q_0,\ldots,Q_{m-1}$ be orthogonal projections,
    and suppose the $m$-block product identity
    \begin{align}
        A_u^{i_1}\cdots A_{u+m-1}^{i_m}
        &=
        c_u\, B_{u+q}^{i_1}\cdots B_{u+q+m-1}^{i_m}
        \notag
    \end{align}
    holds for every base sector $u$ and every length-$m$ word. Then for every
    word $W$ of length $km$ with $k\ge1$,
    $F_u^{W}=c_u^{\,k} \widehat F_{u+q}^{W}$, where $\widehat F$ is the
    repeated corner product built from $B$. Each of the $k$ successive
    length-$m$ blocks is anchored at the same base sector $u$, because a full
    cycle of unit shifts is trivial, $m\cdot 1=0$ in $\mathbb Z_m$; hence every
    block contributes the same scalar $c_u$.
\end{lemma}

\begin{proof}\leanok
    \uses{lem:periodic_corner_product_append}
    Induct on $k$. The base case $k=1$ is the hypothesis. For the inductive
    step split $W$ into its first length-$m$ block and the remaining length-$km$
    word, apply the concatenation law, and use that the junction shift
    $m\cdot 1=0$ returns to the base sector before applying the block identity
    and the induction hypothesis.
\end{proof}

\begin{lemma}[Right-inverse contraction]
    \label{lem:periodic_corner_contraction}
    \lean{MPSTensor.cornerProd_contraction}
    \leanok
    \uses{def:periodic_corner_letter, def:periodic_corner_product}
    Suppose the coefficients $\Omega_{k+1}$ recover each inserted matrix $X_k$
    from the repeated corner product $F_{k+1}$ over length-$L$ words:
    $\sum_{\mathbf j}(\Omega_{k+1}(X_k))_{\mathbf j}\,F_{k+1}^{\mathbf j}=X_k$
    (\cite[eq.~Omegauprop]{DeLasCuevas2017Irreducible} at $X_k$). Then for any
    sector indices $\sigma$,
    \begin{align}
        \sum_{\boldsymbol\rho}
        \left(\prod_k (\Omega_{k+1}(X_k))_{\rho_k}\right)
        \prod_k A_k^{\sigma_k}\,F_{k+1}^{\rho_k}
        &=
        \prod_k A_k^{\sigma_k}\,X_k.
        \notag
    \end{align}
    Summing the concatenated chain against these coefficients collapses each
    repeated product $F_{k+1}$ to its inserted matrix $X_k$. The recovery is
    needed only at the inserted matrices $X_k$, which in the Case-3 contraction
    are the corner-supported tensors that $\Omega_{k+1}$ inverts.
\end{lemma}

\begin{proof}\leanok
    \uses{def:periodic_corner_product}
    Expand each
    $A_k^{\sigma_k}X_k
    =\sum_{\mathbf j}(\Omega_{k+1}(X_k))_{\mathbf j}
    A_k^{\sigma_k}F_{k+1}^{\mathbf j}$ using the right inverse, apply the
    noncommutative distributivity identity
    \begin{align}
        \prod_k\left(\sum_{\mathbf j} f_{k,\mathbf j}\right)
        &=
        \sum_{\boldsymbol\rho}\prod_k f_{k,\rho_k},
        \notag
    \end{align}
    and pull the scalar coefficients
    $\prod_k(\Omega_{k+1}(X_k))_{\rho_k}$ out of each ordered product.
\end{proof}

\begin{lemma}[Inverse cyclic reindexing to the off-diagonal convention]
    \label{lem:periodic_negreindex_shift}
    \lean{MPSTensor.negReindex_paper_shift}
    \leanok
    \uses{thm:offdiag_shift_adjoint_cyclic_shift}
    Let $A$ be left-canonical, $\sum_i A_i^\dagger A_i=1$, with orthogonal sector
    projectors $P_k$ satisfying the adjoint shift $\mathcal E^*(P_{k+1})=P_k$,
    which gives the inverse-indexed grading $P_{k+1}A^i=A^iP_k$. After the
    inverse reindexing $P_k\mapsto P_{-k}$ the projectors satisfy the
    off-diagonal convention $P_{-k}A^i=A^iP_{-(k+1)}$ of
    \cite[eq.~Auprop]{DeLasCuevas2017Irreducible}, under which the
    corner-transition products are nonzero.
\end{lemma}

\begin{proof}\leanok
    Apply the grading at sector $-(k+1)$ and simplify the cyclic index.
\end{proof}

\begin{lemma}[Telescoping of a corner product]
    \label{lem:periodic_cornerprod_telescope}
    \lean{MPSTensor.cornerProd_eq_conj_evalWord}
    \leanok
    \uses{def:periodic_corner_product}
    Under the off-diagonal convention $P_kA^i=A^iP_{k+1}$, with the $P_k$
    orthogonal projections, the repeated corner-transition product collapses
    through its intermediate projectors:
    \begin{align}
        F_u^{\mathbf i}
        &=
        P_u\,A^{\mathbf i}\,P_{u+|\mathbf i|},
        \notag
    \end{align}
    where $A^{\mathbf i}$ is the ordinary word product and the right projector
    records the accumulated sector offset.
\end{lemma}

\begin{proof}\leanok
    \uses{def:periodic_corner_product}
    Induct on the word. The head junction $P_uA^iP_{u+1}=P_uA^i$ collapses by
    the shift and idempotency; the induction hypothesis handles the tail, and
    the accumulated offset realigns.
\end{proof}

\begin{lemma}[Blocked letter as a diagonal corner]
    \label{lem:periodic_cornerprod_blockdiag}
    \lean{MPSTensor.cornerProd_eq_blockDiagCorner}
    \leanok
    \uses{lem:periodic_cornerprod_telescope}
    Under the off-diagonal convention $P_kA^i=A^iP_{k+1}$, with the $P_k$
    orthogonal projections, the corner-transition product of a single blocked
    letter is the diagonal blocked corner of
    \cite[eq.~Cu]{DeLasCuevas2017Irreducible}:
    $F_u^{(I)}=P_u\,(A^{[m]})^I\,P_u$.
    Combined with the corner-isomorphism letter identity, this exhibits the
    compressed sector letter as a corner-transition product.
\end{lemma}

\begin{proof}\leanok
    \uses{lem:periodic_cornerprod_telescope}
    A blocked letter is a word of length exactly $m$; the accumulated offset
    $m\cdot1=0$ vanishes, so the telescoping gives the diagonal corner.
\end{proof}

\begin{definition}[Flattening bijection on word indices]
    \label{def:periodic_block_flatten_equiv}
    \lean{MPSTensor.blockFlattenEquiv}
    \leanok
    The bijection between blocked words of length $L$ over the alphabet of
    $d^m$ blocked letters and single-site words of length $mL$ over the alphabet
    of $d$ letters, $[d^m]^L\simeq[d]^{mL}$, is obtained from the iterated
    blocking bijection. Here $[q]^r$ denotes the set of words of length $r$
    over a $q$-letter alphabet.
\end{definition}

\begin{lemma}[Flattening bijection agrees with word flattening]
    \label{lem:periodic_block_flatten_word_compat}
    \lean{MPSTensor.ofFn_blockFlattenEquiv}
    \leanok
    \uses{def:periodic_block_flatten_equiv}
    If $W\in[d^m]^L$ is a blocked word and $\widehat W\in[d]^{mL}$ is its
    image under the flattening bijection, then $\widehat W$ is exactly the
    concatenation of the $L$ decoded length-$m$ blocks of $W$.
\end{lemma}

\begin{proof}\leanok
    The decoding bijection identifies a word-indexing function with its ordered
    list of letters, and the iterated blocking bijection concatenates the
    decoded length-$m$ blocks in order.
\end{proof}

\begin{lemma}[Corner product of a flattened word]
    \label{lem:periodic_cornerprod_flatten}
    \lean{MPSTensor.cornerProd_flatten_eq_phi_evalWord}
    \leanok
    \uses{def:periodic_corner_product, def:periodic_block_flatten_equiv}
    Let $P_k$ be orthogonal projections satisfying the off-diagonal convention
    $P_kA^i=A^iP_{k+1}$, and let $\varphi_k$ be the multiplicative corner
    isomorphisms with $\varphi_k(\bar A_k^I)=P_k(A^{[m]})^IP_k$. Then for a
    nonempty blocked word $W$, with flattening $\widehat W$, the
    corner-transition product is the corner image of the word evaluation,
    \begin{align}
        F_k^{\widehat W}
        &=
        \varphi_k(\bar A_k^{\,W}),
        \notag
    \end{align}
    the content of \cite[eq.~Fu]{DeLasCuevas2017Irreducible} read through
    $\varphi$.
\end{lemma}

\begin{proof}\leanok
    \uses{lem:periodic_cornerprod_blockdiag}
    Induct on $W$. The head blocked letter collapses by the blocked-letter
    diagonal corner to $\varphi_k(\bar A_k^I)$, and multiplicativity of
    $\varphi$ reassembles the letters:
    \begin{align}
        F_k^{\widehat W}
        &=
        \prod_{I\in W}\varphi_k(\bar A_k^I)
        =
        \varphi_k\left(\prod_{I\in W}\bar A_k^I\right)
        =
        \varphi_k(\bar A_k^{\,W}).
        \notag
    \end{align}
    The empty word is excluded so the induction stays inside the corner.
\end{proof}

\begin{lemma}[Ambient corner right inverse from the sector right inverse]
    \label{lem:periodic_ambient_corner_inverse}
    \lean{MPSTensor.exists_ambientCornerRightInverse_of_sectorRightInverse}
    \leanok
    \uses{lem:periodic_cornerprod_flatten,
        lem:periodic_block_flatten_word_compat}
    From a sector right inverse $\Omega$ recovering every sector word product
    $\bar A_k^{\,W}$ over blocked words $W\in[d^m]^L$ with $L>0$
    (\cite[eq.~Omegauprop]{DeLasCuevas2017Irreducible}) one obtains an ambient
    corner right inverse $\hat\Omega$ over single-site words of length $mL$.
    For every corner-supported $Y$, that is, $P_kYP_k=Y$,
    \begin{align}
        \sum_{\mathbf j\in[d]^{mL}}
            (\hat\Omega_k(Y))_{\mathbf j}\,F_k^{\mathbf j}
        &=
        Y.
        \notag
    \end{align}
\end{lemma}

\begin{proof}\leanok
    \uses{lem:periodic_cornerprod_flatten,
        lem:periodic_block_flatten_word_compat}
    Set $\hat\Omega_k(Y)=\Omega_k(\varphi_k^{-1}(Y))$ for corner-supported $Y$.
    Applying $\varphi_k$ to the sector recovery and using the corner product of a
    flattened word,
    \begin{align}
        \sum_{W}(\hat\Omega_k(Y))_{\widehat W}\,F_k^{\widehat W}
        &=
        \varphi_k\Bigl(        \sum_{W}(\Omega_k(\varphi_k^{-1}(Y)))_{W}
            \bar A_k^{\,W}
        \Bigr)
        =
        \varphi_k(\varphi_k^{-1}(Y))
        =
        Y.
        \notag
    \end{align}
\end{proof}

\begin{lemma}[Sector contraction]
    \label{lem:periodic_sector_contraction}
    \lean{MPSTensor.exists_cornerProd_contraction_of_sectorRightInverse}
    \leanok
    \uses{lem:periodic_ambient_corner_inverse, lem:periodic_corner_contraction}
    Feeding the ambient corner right inverse into the right-inverse contraction
    yields the contracted product-tensor identity
    \cite[eq.~resultprop]{DeLasCuevas2017Irreducible}. For corner-supported
    inserted matrices $X_k$,
    \begin{align}
        \sum_{\boldsymbol\rho}
        \left(\prod_k(\hat\Omega_{k+1}(X_k))_{\rho_k}\right)
        \prod_k A_k^{\sigma_k}F_{k+1}^{\rho_k}
        &=
        \prod_k A_k^{\sigma_k}X_k.
        \notag
    \end{align}
\end{lemma}

\begin{proof}\leanok
    \uses{lem:periodic_ambient_corner_inverse, lem:periodic_corner_contraction}
    Apply the right-inverse contraction with the ambient corner right inverse in
    place of the abstract right inverse.
\end{proof}
